\zlabel{triptych:concise:commentary:start}
\section{The Propers: Detailed Commentary}

\subsection{The command and the Lord who gives it}

The Gospel commands an undivided love, while the Mass asks God to heal those who must live it. Its opening confession of divine justice becomes a plea for mercy; its final prayer seeks the cure of vice and an eternal remedy. Between them, the worshippers hear whose people they are, whom they must love, and what fidelity requires. The command reaches the whole person, but it is addressed within a common life sustained by gifts.

Three readings bring different needs into focus. \emph{Healed charity} asks how the heart becomes able to love what it knows it ought to love. \emph{A gathered people} asks why those receiving the same Lord, faith, and baptism must bear one another. \emph{Faithful offering} asks how confession, prayer, and sacramental reception become continuing fidelity. They are complementary emphases: the person being healed belongs to a body, and that body worships in order to live from what it receives.

Matthew 22:34--46 gives all three their center. Jesus answers a lawyer's test with love of God in the whole heart, soul, and mind, and love of neighbor as oneself. Chrysostom joins the commands through obedience: love of God cannot coexist with refusal of the neighbor whom God commands us to love. Augustine explains their common direction. The good sought for oneself in God is also sought for the neighbor, including an enemy; the neighbor never becomes a disposable means to one's own good.\footnote{Chrysostom, \textit{Matthew}, hom.~71, on 22:37--40; Augustine, \textit{De doctrina christiana}, I.22.20--21; I.30.31--33. Full edition details appear in References.}

The answer may change its questioner. Chrysostom reads the favorable account in Mark as evidence of improvement during the exchange. Augustine allows that possibility but also an initially cautious test rather than malicious deception. Matthew's test does not settle every question about the lawyer's interior disposition. Religious knowledge can become obedience; knowing the answer does not prove that it has.\footnote{Chrysostom, \textit{Matthew}, hom.~71; Augustine, \textit{De consensu evangelistarum}, II.73.141--142.}

Jesus then asks whose son the Christ is. Davidic descent is true, yet David calls him Lord in Psalm 109:1. Chrysostom takes this as disclosure of Christ's divinity without denial of his descent. Healed charity therefore meets its source in the divine Lord who commands it. For the gathered people, he is the one Lord above every faction. For faithful offering, his claim makes worship no payment that leaves the rest of the giver's life self-owned.

The difference between a test and an encounter matters to all three readings. Healed charity asks what the answer does to the questioner: the command must reach the person using it. The gathered people asks how religious knowledge is shared: an answer intended to defeat another can become a truth received in common. Faithful offering asks whether hearing leaves the hearer otherwise self-owned. In each case the Gospel's second exchange presses beyond admiration for an excellent moral reply. David's Lord addresses the one who has been discussing him.

The witnesses still differ concerning motive. Chrysostom narrates improvement; Augustine retains another possible starting point. Both attend to a person's response, and neither reads every person's hidden intention. The command's demand is accordingly exacting, and it falls on the hearer's own will, not on another's hidden mind. A hearer can ask whether he is willing to obey the truth he knows, while leaving another person's initial interior state to God. Chrysostom's and Augustine's alternatives both leave that interior state with God; they are two different stories, not one psychology told twice.\footnote{Chrysostom, \textit{Matthew}, hom.~71; Augustine, \textit{De consensu evangelistarum}, II.73.141--142.}

\subsection{What enables the love that is commanded?}

The Introit joins Psalm 118:137,124 with verse 1: God's judgment is right, his servant needs mercy, and blessedness belongs to those walking his law. Augustine keeps justice and mercy together. The servant does not appeal to righteousness achieved independently of God; he asks for the teaching by which God makes him righteous. Bellarmine gives that teaching a practical edge: it includes love of the command, beyond knowing its content. Healing reaches desire as well as understanding.\footnote{Augustine, \textit{Psalms}, Ps.~118, English paragraphs 123--124,136; Bellarmine, \textit{Psalms}, Ps.~118 ad 124.}

The three readings consequently hear more than one demand in this opening. Healed charity relinquishes self-sufficiency. The gathered people places the singular servant among the blessed walkers: nobody's need exempts him from common life. Faithful offering begins when confession ceases defending what God condemns. Augustine's account of confession in Psalm 75 makes this peace with God concrete. Mercy rescues the sinner; it does not vindicate his wrong.\footnote{Augustine, \textit{Psalms}, Ps.~75, paragraph 3.}

The Collect asks that God's people escape diabolical contagion and follow God alone with a pure mind. Its petition concerns both freedom from corruption and undivided allegiance. Paul's Epistle, Ephesians 4:1--6, shows that allegiance in humility, meekness, patience, and mutual bearing in charity. Chrysostom brings humility down to tone of voice and treatment of enemies. Thomas names pride, anger, impatience, and disordered zeal as threats to peace; bearing another includes correction at a fitting time. Patient charity therefore seeks the other's good, rather than either humiliating him or agreeing negligently with his wrong.\footnote{Chrysostom, \textit{Ephesians}, hom.~9; Thomas, \textit{Super Ephesios}, ch.~4, lect.~1.}

This is the moral meeting point of the readings. The healed heart tells a necessary truth without contempt. The gathered people bears differences without surrendering the good that unites it. Faithful offering carries prayer into an existing responsibility, even after enthusiasm has faded. The conduct is real, while the need for grace remains real. An absence of emotional ease does not make an act of patient love worthless.

The Secret and Postcommunion make the source of that help explicit: the holy action is asked to free the people from offenses, and the mysteries to heal vices. Thomas explains Eucharistic nourishment through Christ present and his Passion: grace is sustained and charity strengthened. His teaching illuminates these petitions. The demand for fidelity remains inside a gift that makes fidelity possible.\footnote{Thomas, \textit{Summa theologiae}, III, q.~79, a.~1, corpus.}

The literal forms of the prayers reinforce the distinction between help and replacement. The Collect asks that the people follow God; the Secret seeks freedom from offenses; the Postcommunion asks that vices be cured. Following, preservation and healing are related, but each names a different need. The Gospel does not become a prayer that somebody else love in the hearer's place. The prayers ask that the hearer become able to live what the command requires. Nor does the command make the petitions redundant once the correct course is known.

That difference gives Bellarmine's practical teaching special force. More information can be necessary, but the person who already knows that contempt is wrong needs a further conversion of desire. Augustine's insistence on continuing illumination likewise refuses the idea that a past understanding makes present dependence unnecessary. Healed charity receives light and strength anew; common life gives that dependence a testing ground; faithful offering carries it into a responsibility that remains after a moment of conviction. The repeated prayer therefore need not mean that an earlier gift failed. It can acknowledge that life remains dependent on its giver.\footnote{Augustine, \textit{Psalms}, Ps.~118, paragraphs 123--124; Bellarmine, \textit{Psalms}, Ps.~118 ad 124.}

\subsection{How can many persons pray as one people?}

Paul's common gifts precede the effort to preserve unity: one body, Spirit, hope, Lord, faith, baptism, and Father. Chrysostom describes the Spirit joining persons separated by wealth, age, and manner of life; Thomas pictures a city with one ruler, law, signs, and end. The first emphasizes the life that joins; the second its common order. Neither makes a congregation's agreement the creator of its Lord.\footnote{Chrysostom, \textit{Ephesians}, hom.~9; Thomas, \textit{Super Ephesios}, ch.~4, lect.~2.}

Chrysostom offers several explanations of the one spirit, including the received Spirit, concord, and zeal. In his account the giver and the people's disposition are distinct: enthusiasm alone cannot establish holiness. Thomas accordingly excludes unity in evil. Chrysostom also distinguishes gifts common to salvation from different gifts for ministry. Equal participation in the common good neither abolishes different offices nor makes a larger task a private trophy.\footnote{Chrysostom, \textit{Ephesians}, homs.~10--11; Thomas, \textit{Super Ephesios}, ch.~4, lect.~1.}

The Gradual, Psalm 32:12,6, places the chosen people beside heavens established by Word and Spirit. Bellarmine begins with created heavens and inseparable Trinitarian action. Augustine identifies spiritual heavens in the righteous, whose firmness comes through Word and Spirit. Creation praise is the literal foundation; Augustine reads the holy people's dependence spiritually beside it, not in its place. Augustine's elected people belongs to the heavenly city, rather than supplying a modern state's claim to superiority.\footnote{Augustine, \textit{Psalms}, Ps.~32, first exposition, paragraphs 6,12; Bellarmine, \textit{Psalms}, Ps.~32 ad 6,12.}

The Alleluia gives this people a singular cry from Psalm 101:2. The psalm begins in affliction and reaches toward Sion's restoration and gathered peoples. Augustine hears Christ with his members: their poverty belongs within his body's voice, while their personal sins are never sins of the Head. Bellarmine stresses that prayer itself needs grace. For healed charity, even the first cry receives help; for the gathered people, one sufferer belongs within common praise; for faithful offering, supplication shapes the listening that follows.\footnote{Augustine, \textit{Psalms}, Ps.~101, paragraphs 1--3; Bellarmine, \textit{Psalms}, Ps.~101, opening commentary.}

Daniel's Offertory carries that solidarity into intercession for sanctuary and people. The chant adapts Daniel 9:17--19; the full chapter confesses communal failure and appeals to mercy. Jerome's commentary on verse 20 allows either confession as one of the people or humble identification without personal participation in the sin confessed. These are different accounts of culpability, though both give solidarity an active voice. Asking mercy for the community neither assigns everyone equal guilt nor permits the righteous to dismiss its wounds.\footnote{Jerome, \textit{Commentaria in Danielem}, II, on 9:17--18,20.}

In the Gradual, too, the witnesses read at two levels, not in a single metaphor. Bellarmine's heavens belong first to creation, whereas Augustine's righteous-as-heavens reading draws a spiritual relation from their dependence on Word and Spirit. The sky need not be secretly a congregation for the spiritual reading to hold. The literal praise gives a real foundation: what exists is not its own source. Augustine then makes that dependence searching for those tempted to take credit for holiness. The elected people can no more manufacture its final good than the created heavens can give themselves existence.\footnote{Augustine, \textit{Psalms}, Ps.~32, first exposition, paragraphs 6,12; Bellarmine, \textit{Psalms}, Ps.~32 ad 6,12.}

Paul's community must preserve a unity that remains larger than a local circle. Chrysostom includes the faithful across places and generations; Thomas makes the shared ruler, faith, sacramental signs and end the city's order. These emphases correct different reductions. A warm feeling among familiar people is too narrow for the body Chrysostom describes. Admiration for a particular leader is too fragile a basis for the common gifts Thomas identifies. Their accounts support a charity that receives real difference in a body that is no collection of unrelated interests.\footnote{Chrysostom, \textit{Ephesians}, hom.~10; Thomas, \textit{Super Ephesios}, ch.~4, lect.~2.}

That common life does not erase the scale of one person's suffering. The Alleluia retains a singular petition; Augustine's Head-and-members reading does not require the sufferer to abandon his own voice. Daniel's intercession then gives common need a named object. The first voice can belong within a body, and the second can represent a people, though neither is every other person's confession. The gathered-people reading therefore includes a distinction that the healed-heart reading needs: shared prayer carries different wounds. The offering reading adds that carrying them cannot stop at observing them. Intercession is an act in which the worshipper accepts another's good as a concern of his own.\footnote{Augustine, \textit{Psalms}, Ps.~101, paragraphs 1--3; Jerome, \textit{Daniel}, II, on 9:17--20.}

\subsection{What must follow an offering?}

The literal speakers remain distinct: the afflicted psalmist, Daniel, Paul, and Jesus' interlocutors do not occupy one historical scene. Yet the Mass's actual sequence joins confession, listening, intercession, holy action, reception, and prayer for lasting fruit. Faithful offering follows that movement most directly. Healed charity asks what changes in the giver along the way; the gathered people asks whose burdens his offering includes.

Jerome reads Daniel's sanctuary petition as a request for promised restoration: divine faithfulness becomes a reason to pray. His reading of the evening sacrifice at verse 21 emphasizes persevering prayer. Augustine, discussing Psalm 75, connects remembrance of mercy with baptism and daily sacrifice. Worship sustains a history of received grace; its repetition calls the healed person to remember the giver through the life that follows.\footnote{Jerome, \textit{Daniel}, II, on 9:17,21; Augustine, \textit{Psalms}, Ps.~75, paragraphs 9--10.}

The Secret's petition concerning future offenses needs a precise distinction. Thomas says that the Eucharist preserves by strengthening spiritual life and resisting hostile powers, while free will can subsequently turn away. Cummiskey's historical English concerning sins yet to be committed therefore supplies no advance absolution or guarantee of perseverance. Healed charity remains a task for a free person; the body's members still need forbearance; the worshipper must live the life nourished by the sacrament.\footnote{Thomas, \textit{Summa}, III, q.~79, a.~6, corpus and ad 1; Cummiskey, \textit{Roman Missal} (1861), Secret for this Sunday, p.~433.}

The Communion, Psalm 75:12--13, commands vows and their payment and gathers gift-bearers around the feared Lord. Augustine distinguishes common Christian fidelity from particular vows: all owe faithful living, while a special undertaking adds its own obligation. God's help enables fulfillment. The vow paid need not be a new private one; it may be an existing duty finally honored. A higher profession cannot make pride holy.\footnote{Augustine, \textit{Psalms}, Ps.~75, paragraph 11.}

Augustine then reads those around the Lord as people who recognize truth as common, refusing to possess it for a faction. Here faithful offering and the gathered people meet: the gift is brought to a shared center that the giver cannot own. Healed charity meets them in the renunciation of pride. The same chant places rulers under judgment. Bellarmine emphasizes God's concrete authority over kings' lives; Augustine applies the humbled spirit of princes to pride and turns rule toward one's own body. Public sovereignty and inward self-rule are different emphases, neither exhausted by the other.\footnote{Augustine, \textit{Psalms}, Ps.~75, paragraphs 12--13; Bellarmine, \textit{Psalms}, Ps.~75 ad 11--12.}

The three readings weigh the same Communion differently. Healed charity hears a test of growth: the cure must become actual fidelity. The gathered people sees gift-bearers around a common Lord: truth is received and shared, not owned as a means of making a faction. Faithful offering hears speech becoming accountable: an undertaking is fulfilled rather than merely repeated. The differences give one another necessary checks. A private account of progress can forget the neighbor; concern for the group can evade personal duty; a conspicuous undertaking can hide pride. The chant brings the giver, the community and the promised act before one judge.

Augustine's distinction between common fidelity and special vows does not call for ever more dramatic promises. A person can owe the difficult good already present in ordinary responsibilities. His comparison of a humble married person and a proud consecrated virgin makes profession answer to charity rather than substituting for it. Bellarmine's attention to the lives of rulers gives the same judgment an outward reach: power cannot purchase permanence. Augustine's inward application to pride does not cancel that public force, and Bellarmine's public emphasis does not release the ordinary believer from self-rule.\footnote{Augustine, \textit{Psalms}, Ps.~75, paragraphs 11--13; Bellarmine, \textit{Psalms}, Ps.~75 ad 11--12.}

\subsection{What is the end of this healing?}

The Postcommunion asks together for cured vices and an eternal remedy. Present improvement is necessary but does not exhaust that hope. Thomas directs sacramental nourishment toward perfect peace and unity in glory; Bellarmine directs the chosen people's faith, hope, and charity toward vision. Judgment remains within the promise: gifts and religious standing cannot purchase exemption from the Lord.\footnote{Thomas, \textit{Summa}, III, q.~79, a.~2, corpus; Bellarmine, \textit{Psalms}, Ps.~32 ad 12.}

The four senses unfold differently within each reading. Literally, healed charity hears commands beside petitions for help; the gathered people hears distinct speakers called into communal life; faithful offering follows acts of confession, supplication, and commitment. Allegorically, their center is Christ: the divine Lord who heals, the Head whose body prays, and the Savior whose Passion gives the mysteries efficacy. Morally, they call respectively for an undivided heart, patient common life, and worship fulfilled in conduct. Anagogically, love's cure, the people's inheritance, and the offering's lasting fruit reach their completion in communion with God and the saints.

These are differences in what the same hope asks of the present. A person learning patience receives healing; a community practicing it keeps a given unity; a worshipper persevering in it honors what he has offered. None becomes self-sufficient. The people that hears the command to love goes on asking for the remedy by which love can endure.
